\documentclass[11pt,a4paper]{article}
\usepackage[margin=0.86in]{geometry}
\usepackage{amsmath,amssymb,booktabs,array,enumitem,fancyhdr,hyperref}
\usepackage[T1]{fontenc}
\usepackage{lmodern}
\usepackage{microtype}
\usepackage{xcolor}
\hypersetup{colorlinks=true,urlcolor=blue!55!black,linkcolor=black}
\setlength{\parindent}{0pt}
\setlength{\parskip}{5pt}
\setlist{itemsep=1pt,topsep=3pt}
\pagestyle{fancy}
\fancyhf{}
\fancyhead[L]{Homework 8}
\fancyhead[C]{CIS 5511}
\fancyhead[R]{Md. Mizanur Rahman}
\fancyfoot[C]{\thepage}
\renewcommand{\headrulewidth}{0.4pt}
\newcommand{\code}[1]{\texttt{#1}}

\begin{document}
\thispagestyle{plain}
\begin{center}
  {\large CIS 5511 -- Homework 8}\\[3pt]
  Md. Mizanur Rahman\\[-1pt]
  \texttt{md.mizanur.rahman@temple.edu}
\end{center}
\vspace{6pt}
{\large\bfseries Longest Common Substring Using Dynamic Programming}\par

\section*{1. Data structure and recursive formulation}
Let $X=x_1x_2\cdots x_n$ and $Y=y_1y_2\cdots y_m$ be two strings.
A \emph{substring} is a contiguous block of characters. The dynamic programming
table $D$ has $(n+1)$ rows and $(m+1)$ columns, where
\[
 D[i,j]=\text{length of the longest common \emph{suffix} of }
 X[1\ldots i]\text{ and }Y[1\ldots j].
\]
This suffix-based state enforces contiguity: only matches that end at positions
$i$ and $j$ can be extended. The recurrence is
\[
 D[i,j]=\begin{cases}
 0, & i=0\text{ or }j=0,\\
 D[i-1,j-1]+1, & i,j>0\text{ and }x_i=y_j,\\
 0, & i,j>0\text{ and }x_i\ne y_j.
 \end{cases}
\]
The solution length is $L=\max_{0\leq i\leq n,\,0\leq j\leq m}D[i,j]$.
For a maximizing cell $(i^*,j^*)$, the answer is the $L$ characters ending at
$i^*$ in $X$ (equivalently, at $j^*$ in $Y$). The implementation saves $L$
and the two zero-based start indices. For multiple optimal substrings it
selects the earliest start in $X$, breaking further ties by earliest start in
$Y$. If $L=0$, the substring is empty and both positions are \code{None}.

\textbf{Difference from the Longest Common Subsequence (LCS).}
Both problems compare two prefixes and reuse solutions to smaller subproblems
in a two-dimensional DP table. However, an LCS may skip characters, while an
LCStr must use consecutive characters in both strings. The usual LCS state
$C[i,j]$ is the best \emph{subsequence length anywhere} in the two prefixes:
\[
 C[i,j]=\begin{cases}
 0, & i=0\text{ or }j=0,\\
 C[i-1,j-1]+1, & x_i=y_j,\\
 \max\{C[i-1,j],C[i,j-1]\}, & x_i\ne y_j.
 \end{cases}
\]
Thus LCS carries forward the best answer when characters differ and returns
$C[n,m]$. LCStr resets a mismatched cell to zero and returns the \emph{maximum
over all cells} of $D$, not necessarily $D[n,m]$. For example,
$X=\code{ABCDEF}$ and $Y=\code{ACE}$ have LCS \code{ACE} (length $3$),
but LCStr \code{A} (length $1$; \code{C} and \code{E} are also optimal).

\section*{2. Dynamic programming algorithm}
The implemented \code{longest\_common\_substring(x,y)} uses the following
bottom-up pseudocode in the textbook's numbered-line format. Characters of
$X$ and $Y$ are indexed from $1$; row $0$ and column $0$ of $D$ represent
empty prefixes. The returned start positions are zero-based to agree with
the Python implementation and the submitted testing output.

\begingroup\small
\begin{verbatim}
LONGEST-COMMON-SUBSTRING(X, Y)

 1  n = X.length
 2  m = Y.length
 3  D = new array[0:n, 0:m]

 4  for i = 0 to n
 5      D[i, 0] = 0
 6  for j = 0 to m
 7      D[0, j] = 0

 8  maxLength = 0
 9  endX = 0
10  endY = 0

11  for i = 1 to n
12      for j = 1 to m
13          if X[i] == Y[j]
14              D[i, j] = D[i-1, j-1] + 1
15              if D[i, j] > maxLength
16                  maxLength = D[i, j]
17                  endX = i
18                  endY = j
19          else
20              D[i, j] = 0

21  if maxLength == 0
22      return ("", 0, NIL, NIL, D)

23  startX = endX - maxLength
24  startY = endY - maxLength
25  substring = X[startX+1:endX]
26  return (substring, maxLength, startX, startY, D)
\end{verbatim}
\endgroup

The assignment symbol is \code{=}, while equality is tested with
\code{==}. The subarray notation on line 25 includes both endpoints.
The table is traversed in increasing $i$, then increasing $j$. Since lines
15--18 update the answer only when a \emph{strictly longer} substring is found,
ties retain the earliest starting position in $X$, followed by the earliest
starting position in $Y$. This agrees with the Python program's tie-breaking.

\textbf{Correctness.} The zero row/column are correct for empty prefixes.
For $i,j>0$, unequal final characters cannot terminate a common nonempty
substring, so the correct suffix length is $0$. Equal final characters extend
the longest common suffix of the immediately preceding prefixes by one.
Induction over the table therefore establishes every $D[i,j]$; taking its
maximum finds the globally longest common substring.

\section*{3. Time and space analysis}
Initialization uses $(n+1)(m+1)$ zero entries. Each of the $nm$ inner-loop
iterations performs a comparison, a constant number of table accesses, and
constant-time best-position updates. Therefore the DP computation takes
$\Theta((n+1)(m+1))$ time and space, customarily stated as
$\boxed{\Theta(nm)}$ for nonempty strings. Extracting the returned substring
costs at most $O(\min(n,m))$ time and space, which does not change the overall
bound. The full table is retained because the assignment asks for its contents.
If only the result were required, two rows (or one carefully updated row) could
reduce DP workspace to $O(\min(n,m))$.

\section*{4. Implementation and testing results}
The Python implementation is submitted as \code{lcstr.py} and uses no external
packages. Run \code{python lcstr.py} to execute all seven cases and print the
input, the result, the two starting positions, and \emph{every cell} of the DP
table. The exact console transcript is supplied separately in
\code{lcstr\_output.txt}; \code{python lcstr.py --save} regenerates it.
Each case also checks the DP recurrence and the reconstructed answer.

\begin{center}
\small
\begin{tabular}{@{}c l l l c l@{}}
\toprule
Case & $X$ & $Y$ & Chosen LCStr & Length & Starts $(X,Y)$\\
\midrule
1 & \code{ABABC} & \code{BABCA} & \code{BABC} & 4 & $(1,0)$\\
2 & \code{ABCDEF} & \code{ACE} & \code{A} & 1 & $(0,0)$\\
3 & \code{BANANA} & \code{ANANAS} & \code{ANANA} & 5 & $(1,0)$\\
4 & \code{ABABA} & \code{BABAB} & \code{ABAB} & 4 & $(0,1)$\\
5 & \code{ABC} & \code{XYZ} & empty & 0 & $(-,-)$\\
6 & empty & \code{DATA} & empty & 0 & $(-,-)$\\
7 & \code{NETWORK} & \code{NETWORK} & \code{NETWORK} & 7 & $(0,0)$\\
\bottomrule
\end{tabular}
\end{center}

For illustration, the complete table for Case 2 is shown below. Every
nonzero entry equals $1$ because no two matching characters are adjacent in
both strings; the longest common \emph{subsequence}, in contrast, has length $3$.
\begin{center}
\small
\renewcommand{\arraystretch}{1.05}
\begin{tabular}{c|rrrr}
 $D[i,j]$ & $-$ & A & C & E\\\hline
 $-$ & 0 & 0 & 0 & 0\\
 A & 0 & 1 & 0 & 0\\
 B & 0 & 0 & 0 & 0\\
 C & 0 & 0 & 1 & 0\\
 D & 0 & 0 & 0 & 0\\
 E & 0 & 0 & 0 & 1\\
 F & 0 & 0 & 0 & 0
\end{tabular}
\end{center}
All seven test cases passed. The separate output file contains the complete
$(|X|+1)\times(|Y|+1)$ table for \emph{every} case, including the empty-input
case and ties between multiple longest substrings.

\section*{Files and acknowledgment}
The submission consists of three files:
\code{PT\_08.pdf} (this report), \code{lcstr.py} (source code), and
\code{lcstr\_output.txt} (complete inputs, outputs, and DP tables).
The design follows the dynamic programming and LCS concepts in
\href{https://cis.temple.edu/~pwang/5511-PT/Lecture/5511-10.htm}{Professor Wang's Lecture 10}.

\textbf{AI Use Acknowledgment.} Generative AI (ChatGPT) assisted with Python
implementation, testing, and LaTeX report preparation.
\end{document}
